\chapter{Conclusiones y trabajo futuro} \label{cap5}

\vspace*{5mm}

%F1 Client 1

En esta sección comentaremos las conclusiones a las que llegamos a través de ese análisis y propondremos el trabajo futuro en esta línea.

\section{Conclusiones}

En este trabajo hemos realizado un análisis del estado del arte en el ámbito del aprendizaje federado y en el de la clasificación de imágenes multi-etiqueta poniendo especial atención en el estudio de la heterogeneidad dentro de estos sistemas. Se ha propuesto un modelo para aprendizaje federado multi-etiqueta en el que los clientes tiene diferentes etiquetas usando la aproximación de crear un modelo mixto con capas personalizadas en los clientes y capas compartidas con el servidor. Se ha realizado un estudio experimental con un conjunto de datos multi-etiqueta para comparar el rendimiento con respecto a otro sistema sin federar que tiene la misma configuración. Por último, se han definido una serie de métricas que nos permite evaluar el rendimiento de nuestro experimento.

Nuestra hipótesis inicial era que para la tarea de clasificación de imágenes multi-etiqueta donde varios clientes deben entrenar con un conjunto de datos con etiquetas diferentes, el pertenecer a un entorno federado podría mejorar el rendimiento ya que los distintos clientes necesitan de un aprendizaje común (reconocimiento de bordes, texturas, sombras, etcétera...) que se delega a la parte compartida del modelo mientras la tarea de clasificación en sí permanece en las capas personalizadas de cada cliente. Aplicar un algoritmo como FedPer nos debería permitir atacar el problema del desbalance de etiquetas en los clientes de manera satisfactoria. Al analizar los resultados del experimento observamos que existen dificultades en el entrenamiento del modelo derivados probablemente de la acción de agregación por parte del servidor que afecta negativamente a la convergencia en el cliente. 

A pesar de los resultados negativos, vemos que hay una tendencia ascendente en la mejora del rendimiento conforme avanzan las rondas de servidor por lo que deducimos que aún existe margen de mejora para este experimento, a continuación se enumeran varias opciones:

\begin{description}
\item[1. Algoritmo de agregación] Como hemos comentado anteriormente la diferencia principal entre los dos casos que hemos comparado (Federado vs no federado) es que realizamos la evaluación de cada cliente después de realizar la agregación en el servidor de sus respectivos modelos usando FedAvg. Existen algoritmos que estudian cómo mitigar los efectos de los datos heterogéneos en el momento de la agregación como los algoritmos SCAFFOLD ~\cite{scaffold} y FedProx ~\cite{fedprox}. El uso de algoritmos diferentes podrían mejorar la convergencia de los modelos locales en cada cliente.

\item[2. Épocas locales] Teorizamos que un mayor número de épocas locales afecta a la variación del modelo compartido en mayor medida, minimizando el número de épocas podríamos hacer que la variación de los pesos fuese menor por lo que el impacto que tuviese el entrenamiento en un cliente en otro también disminuiría.

\item[3. Número de rondas federadas] En todos los clientes observamos una mejora conforme avanzan las rondas federadas. Añadiendo más rondas podríamos llegar a un punto en el que los resultados igualasen o mejorasen al caso no federado que se estabiliza en torno al ~0.85 en la métrica F1. Se podrían multiplicar el número de pasos de entrenamiento para terminar de evaluar las limitaciones de este algoritmo.

\item[4. Conjunto de datos] En nuestro estudio pretendemos crear un sistema en el que cada cliente entrene etiquetas distintas, sin embargo, en el conjunto de datos que hemos usado para realizar el experimento también existe un importante desbalance en el número de ocurrencias de cada etiqueta. Eligiendo un conjunto de imágenes multi-etiqueta que no presentase esta característica se obtendrían mejores resultados y facilitaría el análisis evitando la dificultad añadida que nos aporta este desbalance.

\end{description}

\section{Trabajo futuro}

Los datos experimentales obtenidos nos dejan ver que aunque el rendimiento sea peor dentro de un entorno federado que en uno centralizado, para algunas etiquetas el resultado ha sido mejor, por lo que se intuye que hay margen de mejora y aún hay camino que recorrer en esta dirección. Como hemos comentado en el punto anterior tenemos varias opciones para mejorar nuestro modelo. Un sistema que permita a un cliente entrenar una tarea de clasificación en cooperación con otros nodos tendría muchas aplicaciones en el mundo real en distintos campos donde es un requisito la privacidad de los datos como puede ser el procesamiento de imágenes médicas, video-vigilancia u otros ámbitos donde se pueda aplicar la clasificación de imágenes de manera que los costes del entrenamiento puedan ser distribuidos entre los distintos actores.

Una vez asentada la base de nuestro modelo, como podemos observar en las conclusiones, tenemos como tarea pendiente mejorar el rendimiento poniendo el foco en la agregación del servidor, los parámetros de configuración o los datos de entrenamiento. 

Y una vez solventado el problema del rendimiento podemos proponer otro experimento como el de introducir un número de etiquetas distintas para cada cliente. Al delegar la parte de la salida del modelo a las capas personalizadas no debería afectar al entrenamiento global, esto nos daría más flexibilidad a la hora de introducir nuevos clientes con diferentes necesidades o formatos de clasificación.

Por último proponemos la introducción de un cliente sin entrenar en un sistema federado ya entrenado. De esta manera el cliente podría hacer uso de la parte del modelo compartido con el servidor que estaría pre-entrenada y debería de mejorar el entrenamiento de manera más rápida que si entrenase solo ya que solo tendría que ajustar las capas personalizadas en lugar del modelo completo.